% Abstract: limit 200 words and no citations

% currently 198 words
\begin{abstract}
Neuromorphic computing shows promise for advancing computing efficiency and capabilities of AI applications using brain-inspired principles. However, the neuromorphic research field currently lacks standardized benchmarks, making it difficult to accurately measure technological advancements, compare performance with conventional methods, and identify promising future research directions. Prior neuromorphic computing benchmark efforts have not seen widespread adoption due to a lack of inclusive, actionable, and iterative benchmark design and guidelines. To address these shortcomings, we present NeuroBench: a benchmark framework for neuromorphic computing algorithms and systems. NeuroBench is a collaboratively-designed effort from an open community of researchers across industry and academia, aiming to provide a representative structure for standardizing the evaluation of neuromorphic approaches. The NeuroBench framework introduces a common set of tools and systematic methodology for inclusive benchmark measurement, delivering an objective reference framework for quantifying neuromorphic approaches in both hardware-independent (algorithm track) and hardware-dependent (system track) settings. In this article, we outline tasks and guidelines for benchmarks across multiple application domains, and present initial performance baselines across neuromorphic and conventional approaches for both benchmark tracks.
NeuroBench is intended to continually expand its benchmarks and features to foster and track the progress made by the research community.

% % 100 word, more general abstract
% Neuromorphic computing shows promise for advancing computing efficiency and capabilities of AI applications using brain-inspired principles. However, the neuromorphic research field currently lacks standardized benchmarks, making it difficult to accurately measure technological advancements, compare performance with conventional methods, and identify promising future research directions. This article presents NeuroBench, a benchmark framework for neuromorphic algorithms and systems, which is collaboratively designed from an open community of nearly 100 co-authors across over 50 institutions in industry and academia. The NeuroBench framework introduces a common set of tools and systematic methodology for inclusive benchmark measurement, delivering an objective reference framework for quantifying neuromorphic approaches in both hardware-independent and hardware-dependent settings.

% JY: We have hit the word limit above already, no space to expand sentences
% Neuromorphic computing shows promise for advancing computing efficiency and capabilities of AI applications using brain-inspired principles. However, the neuromorphic research field currently lacks standardized benchmarks. This makes it difficult to accurately measure technological advancements, compare performance with conventional methods, and identify promising future research directions. Prior neuromorphic computing benchmark efforts have not seen widespread adoption due to a lack of inclusive, actionable, and iterative benchmark design and guidelines. To address these challenges, we brought together a collaborative community from over 50 institutions in industry and academia. This expert team includes nearly 100 co-authors from relevant domains to develop NeuroBench: a benchmark framework for neuromorphic computing algorithms and systems. In this paper, we collectively present NeuroBench---a benchmark framework to provide a representative structure for standardizing the evaluation of neuromorphic computing approaches, across both algorithms and system architectures. The NeuroBench framework introduces a common set of tools and systematic methodology for inclusive benchmark measurement. It offers both hardware-independent (algorithm track) and hardware-dependent (system track) settings. NeuroBench delivers an objective reference framework for quantifying advancements in neuromorphic computing. We present initial performance baselines across various model architectures on the algorithm track, as well as outline the system track tasks. NeuroBench is intended to continually expand its benchmarks and features to foster and track progress made by the research community.
% However, neuromorphic research lacks standardized benchmarks, making progress difficult to measure and compare. 
% Past benchmarking efforts have failed from lack of inclusive design and adoption.
% This makes it difficult to accurately measure technological advancements, compare current performance to conventional methods, and determine future research directions.
% To address these shortcomings, we present NeuroBench, a collaborative benchmarking framework for neuromorphic algorithms and systems. NeuroBench provides a systematic methodology and common set of tools to evaluate neuromorphic approaches, split into algorithm and hardware tracks.
% Developed through an open community, NeuroBench establishes representative benchmarks and guidelines for measuring neuromorphic progress. The framework is designed to be iterative and extensible, expanding with community input over time.
% NeuroBench delivers a critical foundation and tools to align and accelerate neuromorphic computing advancements. By promoting benchmark standardization, NeuroBench aims to foster breakthroughs in this promising computing paradigm.

% Neuromorphic computing holds great promise in advancing computing efficiency and capabilities of AI applications using brain-inspired principles. 
% However, the rich diversity of methods used in neuromorphic research has led to a lack of standardized benchmarks.
% This makes it difficult to accurately measure technological advancements, compare current performance to conventional methods, and determine future research directions.
% Prior neuromorphic computing benchmark efforts have not seen widespread adoption due to a lack of inclusive, actionable, and iterative benchmark design and guidelines.
% To address these shortcomings, we present NeuroBench: a benchmark framework for neuromorphic computing algorithms and systems. 
% NeuroBench is a collaboratively-designed effort from an open community of academic and industry researchers aiming to provide a representative structure for standardizing the evaluation of neuromorphic approaches.
% The NeuroBench framework introduces a common set of tools and a systematic methodology for inclusive benchmark measurement, split between hardware-independent (\textit{algorithm track}) and hardware-dependent (\textit{system track}) settings. 
% As a whole, the NeuroBench initiative is intended to continually expand its benchmarks and features to foster and align with the progress made by the research community. 
\end{abstract}

